\begin{table}[H]
\centering
\footnotesize
\caption{Tabla 5. Factor de inflación de la varianza de los regresores}
\label{tab:tabla5_vif}
\resizebox{\textwidth}{!}{%
\begin{tabular}{lcc}
\toprule
\textbf{Regresor} & \textbf{VIF} & \textbf{Diagnóstico} \\
\midrule
S\&P 500 & 2.862 & Sin multicolinealidad grave \\
Oro & 1.258 & Sin multicolinealidad grave \\
IGBVL & 1.402 & Sin multicolinealidad grave \\
TC PEN/USD & 1.226 & Sin multicolinealidad grave \\
Var. VIX & 2.639 & Sin multicolinealidad grave \\
Var. UST10 (pbs) & 1.132 & Sin multicolinealidad grave \\
Índice dólar & 1.545 & Sin multicolinealidad grave \\
\bottomrule
\end{tabular}}
\par\vspace{3pt}\parbox{\textwidth}{\scriptsize \textit{Nota.} Criterio: VIF < 10 indica ausencia de multicolinealidad grave. Fuente: elaboración propia con datos de Yahoo Finance, Coin Metrics, BCRPData (BCRP) y FRED (Fed de St. Louis), 2018-2025.}
\end{table}